%% USPSC-Cap4-Metodologia.tex
% ----------------------------------------------------------
% Metodologia da Avaliação
% ----------------------------------------------------------
\chapter{Metodologia da Avaliação}
\label{cap:metodologia}

\section{Considerações Iniciais}
\label{sec:met_consideracoes}

Este capítulo detalha o desenho metodológico da avaliação de experiência do
usuário do aplicativo \textit{Trilha das Árvores}. Descreve o delineamento geral e
seu encadeamento em rodadas, a caracterização do público-alvo por proto-personas,
os critérios de recrutamento e os aspectos éticos, o aparato técnico empregado, o
roteiro de tarefas, o protocolo de condução da observação, os instrumentos de
coleta e o plano de análise.

\section{Delineamento do Estudo}
\label{sec:delineamento}

A avaliação adota delineamento de métodos mistos, de caráter \textbf{formativo},
executado em fases sequenciais e complementares. Conforme delimitado na
Seção~\ref{sec:escopo}, não se pretende produzir medida somativa generalizável,
mas identificar problemas e orientar melhorias verificáveis.

O estudo organiza-se em duas rodadas de avaliação separadas por um ciclo de
correções, segundo o modelo de projeto iterativo:

\begin{alineas}
  \item \textbf{Fase 0 --- Avaliação Heurística.} Inspeção da interface à luz das
  dez heurísticas apresentadas no Quadro~\ref{quad:heuristicas}, com registro de
  cada violação, heurística associada e grau de severidade. As correções
  decorrentes são aplicadas antes do contato com usuários, de modo que o
  participante avalie o conceito de mediação e não defeitos elementares de
  interface.

  \item \textbf{Fase 1 --- Avaliação inicial com usuários.} Estudo observacional
  em campo, sob protocolo \textit{Think-Aloud} híbrido segmentado, seguido de
  entrevista semiestruturada e da aplicação do instrumento TAM.

  \item \textbf{Fase 2 --- Ciclo de correções.} Implementação das melhorias
  priorizadas a partir dos achados da Fase 1, com registro da rastreabilidade
  entre cada correção e o achado que a motivou.

  \item \textbf{Fase 3 --- Reavaliação.} Repetição do protocolo da Fase 1, com
  participantes distintos e sob condições equivalentes, para verificar se os
  atritos corrigidos deixaram de se manifestar.
\end{alineas}

\section{Caracterização do Público-Alvo}
\label{sec:personas}

\subsection{Opção por proto-personas}

Conforme fundamentado na Seção~\ref{sec:personas_teoria}, este trabalho não dispõe
de pesquisa prévia sobre o perfil de visitação do parque. Adotam-se, portanto,
\textbf{proto-personas}: perfis que explicitam as premissas assumidas sobre o
público-alvo, formulados como hipótese a ser verificada com os dados coletados.
Os perfis foram construídos a partir de quatro eixos de diferenciação com
consequência direta sobre a interface: a motivação da visita, a familiaridade com
o parque, a relação com o esforço físico e o perfil de uso do dispositivo em
campo.

% ============================================================================
% >>> [OPCIONAL --- FIGURAS DAS PROTO-PERSONAS] <<<
% Se optar por apresentar as proto-personas em formato de cartão visual,
% gerar as imagens, salvar em USPSC-img/ e incluí-las aqui.
% ============================================================================

\subsection{Proto-persona 1 --- motivação botânica}

Estudante de graduação em área correlata às ciências biológicas, com idade
próxima a 21 anos, externa à ESALQ, em primeira visita ao parque. Sua motivação é
o reconhecimento das espécies: deseja identificar as árvores e obter informação
sobre elas. Não conhece a topografia do campus. Utiliza o celular com desenvoltura,
mas com restrição de bateria e de franquia de dados. O aplicativo lhe é útil na
medida em que identifica corretamente o exemplar diante do qual ela se encontra e
entrega a informação no momento em que ela está parada diante dele. O atrito
esperado decorre da ênfase da interface em métricas de percurso, em detrimento do
conteúdo botânico.

\subsection{Proto-persona 2 --- motivação de esforço físico}

Pós-graduando com idade próxima a 24 anos, praticante regular de corrida,
familiarizado com aplicativos de registro de atividade física. Interpreta a trilha
como percurso cronometrado e demanda fricção mínima: não deseja interromper o
deslocamento para ler a tela. O aplicativo lhe é útil na medida em que suas
métricas sejam legíveis de relance, com o dispositivo em movimento e sob
luminosidade elevada. O atrito esperado decorre de alvos de toque reduzidos, de
contraste insuficiente sob sol direto e da ambiguidade entre distância percorrida
e distância restante.

\subsection{Proto-persona 3 --- motivação de lazer}

Moradora da região, com idade próxima a 26 anos, sem formação na área ambiental,
em visita de fim de semana acompanhada de familiar. Desconhece nomenclatura
científica e não dispõe de disposição para configurações. Conhece o parque
visualmente, mas nunca realizou percurso guiado. Utiliza dispositivo de
especificação intermediária e manifesta reserva quanto à concessão de permissões.
O aplicativo lhe é útil na medida em que a proposta seja compreensível em poucos
segundos e o início não exija auxílio. O atrito esperado decorre da solicitação de
permissões sem justificativa contextualizada e da exibição de dados não
integrados, que compromete a confiança no sistema.

\subsection{Antipersona}

Estudante ou servidor da própria ESALQ, com circulação diária pelo parque e
capacidade de nomear os exemplares sem auxílio. Esse perfil não constitui público
do aplicativo, cuja proposta de valor é a mediação para quem não dispõe desse
repertório. A explicitação da antipersona cumpre função metodológica: fundamenta o
critério de exclusão adotado no recrutamento e responde antecipadamente à
objeção quanto à não utilização de participantes disponíveis no próprio campus.

\section{Participantes, Recrutamento e Aspectos Éticos}
\label{sec:participantes}

\subsection{Dimensionamento da amostra}

A etapa observacional emprega número reduzido de participantes, coerente com a
recomendação consolidada por \citeonline{nielsenlandauer1993} e
\citeonline{virzi1992} de que um pequeno conjunto de usuários revela a maior parte
dos problemas de interface, com ganho marginal decrescente a partir de cinco
participantes. Adotam-se \textbf{três participantes por rodada}, dimensionamento
compatível com o caráter formativo do estudo e com os recursos disponíveis, e
consonante com a orientação metodológica registrada para o projeto.

\subsection{Critérios de inclusão e exclusão}

Os critérios apresentados no Quadro~\ref{quad:recrutamento} derivam diretamente
das proto-personas e da antipersona.

\begin{quadro}[htb]
\centering
\caption{Critérios de recrutamento de participantes}
\label{quad:recrutamento}
\begin{tabular}{p{2.8cm} p{10.0cm}}
\toprule
\textbf{Critério} & \textbf{Definição} \\
\midrule
Inclusão & Idade entre 20 e 26 anos; nunca ter utilizado o aplicativo; nunca ter
realizado percurso guiado no parque; ausência de limitação de mobilidade que
impeça a caminhada; consentimento formal \\
\addlinespace
Exclusão & Vínculo de estudo ou trabalho com a ESALQ; participação no
desenvolvimento do sistema; familiaridade com o parque suficiente para nomear os
exemplares \\
\addlinespace
Variação buscada & Um participante por perfil de motivação, cobrindo as três
proto-personas \\
\addlinespace
Controle & Ausência de familiaridade com o parque mantida constante entre os
participantes, de modo que dificuldades de orientação sejam atribuíveis à
interface ou ao ambiente, e não ao conhecimento prévio do percurso \\
\bottomrule
\end{tabular}
\\[4pt]
\begin{flushleft}
Fonte: Elaborado pelo autor.
\end{flushleft}
\end{quadro}

\subsection{Não reutilização de participantes entre rodadas}

A Fase 3 emprega \textbf{participantes distintos} dos da Fase 1. A justificativa é
o efeito de aprendizagem: parcela substancial do que se avalia é a
aprendibilidade da interface, atributo que existe apenas no primeiro contato. Um
participante que já percorreu a trilha conhece a mecânica de validação por QR
Code, a localização dos controles e o significado dos indicadores; seu eventual
melhor desempenho na segunda rodada seria indistinguível entre efeito do
redesenho e efeito da experiência prévia, confundimento que não admite correção
estatística em amostra desta dimensão. Mantém-se constante, entre as rodadas, o
\textbf{perfil} dos participantes e não sua identidade, de modo que a diferença
observada seja atribuível ao redesenho.

Registra-se a possibilidade de retorno de um participante da primeira rodada, após
a conclusão das sessões da segunda, exclusivamente como verificação confirmatória
pontual quanto à resolução de um problema específico por ele vivenciado. Tal
participação, se ocorrer, será relatada à parte e \textbf{não integrará} a amostra
da segunda rodada.

\subsection{Aspectos éticos}

A participação é voluntária e precedida da assinatura do Termo de Consentimento
Livre e Esclarecido, que assegura anonimato, confidencialidade, uso estritamente
acadêmico dos dados, informação sobre a gravação de áudio e de tela, e direito de
desistência a qualquer momento sem necessidade de justificativa. Não há gravação
de imagem do participante. Os registros são identificados por código sequencial e
armazenados em meio de acesso restrito.

% ============================================================================
% >>> [PREENCHER --- SITUAÇÃO JUNTO AO COMITÊ DE ÉTICA] <<<
% Redigir aqui o desfecho da consulta à coordenação da pesquisa quanto à
% exigência de submissão do protocolo ao Comitê de Ética em Pesquisa (CEP),
% por se tratar de estudo com seres humanos com registro de áudio.
% Informar: se houve submissão, o número do parecer e a data de aprovação;
% ou a justificativa formal de dispensa.
% ============================================================================

\section{Aparato Técnico e Padronização das Condições}
\label{sec:aparato}

\subsection{Padronização do dispositivo}

Todas as sessões utilizam um \textbf{único dispositivo}, fornecido pelo
pesquisador e idêntico para todos os participantes. A decisão elimina a variação
introduzida por diferenças de sistema operacional, dimensão e densidade de tela,
conjunto de ícones do sistema e desempenho de hardware --- variáveis que afetariam
diretamente a usabilidade observada e cuja influência não seria separável dos
efeitos da interface avaliada.

Antes de cada sessão, o aplicativo tem seus dados reinicializados e suas permissões
revogadas, de modo que cada participante encontre o sistema em estado idêntico e
seja exposto aos diálogos de solicitação de permissão, cuja recepção constitui
objeto de observação.

% ============================================================================
% >>> [PREENCHER --- ESPECIFICAÇÃO DOS COMPONENTES TÉCNICOS] <<<
%
% Espaço reservado, conforme planejado, para a descrição do aparato.
% Preencher o Quadro abaixo com os dados reais e redigir, em seguida, um
% parágrafo justificando as escolhas.
%
% Itens a documentar:
%   - Dispositivo: fabricante, modelo, versão do sistema operacional,
%     dimensão e resolução de tela, memória, capacidade de bateria.
%   - Aplicativo do sistema sob teste: número da versão / identificador do
%     build utilizado nas sessões (essencial: a Fase 3 usa outro build, e a
%     comparação entre rodadas exige que ambos estejam identificados).
%   - Gravação de tela: aplicativo utilizado, resolução, taxa de quadros,
%     se captura áudio interno.
%   - Gravação de áudio: modelo do microfone (lapela / direcional),
%     dispositivo de registro, taxa de amostragem, formato do arquivo.
%   - Cronometragem: instrumento independente utilizado.
%   - Conectividade: operadora, tecnologia de acesso disponível no local.
%   - Energia: bateria externa, capacidade.
%   - Material impresso: relação dos formulários levados a campo.
% ============================================================================

\begin{quadro}[htb]
\centering
\caption{Aparato técnico utilizado nas sessões}
\label{quad:aparato}
\begin{tabular}{p{4.2cm} p{8.6cm}}
\toprule
\textbf{Componente} & \textbf{Especificação} \\
\midrule
Dispositivo móvel & \textit{[preencher]} \\
Sistema operacional & \textit{[preencher]} \\
Versão do aplicativo & \textit{[preencher --- Fase 1]} \\
Gravação de tela & \textit{[preencher]} \\
Microfone & \textit{[preencher]} \\
Dispositivo de gravação de áudio & \textit{[preencher]} \\
Cronômetro & \textit{[preencher]} \\
Conectividade & \textit{[preencher]} \\
Alimentação auxiliar & \textit{[preencher]} \\
\bottomrule
\end{tabular}
\\[4pt]
\begin{flushleft}
Fonte: Elaborado pelo autor.
\end{flushleft}
\end{quadro}

\subsection{Padronização das condições ambientais}

Conforme argumentado na Seção~\ref{sec:contexto_movel}, as condições ambientais
integram o contexto de uso e, portanto, o objeto da avaliação. A estratégia
adotada não é sua eliminação, mas sua \textbf{padronização e registro}. Assim:

\begin{alineas}
  \item todas as sessões ocorrem no mesmo período do dia, na mesma trilha, com o
  mesmo ponto de partida e a mesma ordem de tarefas;

  \item são suspensas as sessões em condição de chuva, que inviabiliza a operação
  da tela sensível ao toque e introduz risco ao participante;

  \item as condições efetivas de cada sessão --- data, horário de início e
  término, temperatura, cobertura de nuvens e ocorrência de chuva nas horas
  precedentes --- são registradas em cabeçalho da ficha de observação;

  \item as mesmas condições são replicadas na Fase 3, de modo que a comparação
  entre rodadas não seja confundida por variação ambiental.
\end{alineas}

\subsection{Aferição do esforço físico percebido}

O condicionamento físico dos participantes não é passível de equalização. Adota-se,
por essa razão, sua conversão em \textbf{variável registrada}: ao término das
tarefas, antes da entrevista, aplica-se item único de esforço percebido, adaptado
da escala de percepção subjetiva de esforço de \citeonline{borg1982}, em amplitude
de 0 a 10. O registro permite verificar a homogeneidade do esforço entre
participantes e sustentar, ou refutar, a atribuição dos atritos observados à
fadiga.

\section{Roteiro de Tarefas}
\label{sec:tarefas}

As tarefas são enunciadas na forma de \textbf{cenário}, e não como instrução
contendo o nome do controle a ser acionado. A formulação evita que o enunciado
resolva antecipadamente o problema de localização e reconhecimento que se pretende
observar. O Quadro~\ref{quad:tarefas} apresenta o conjunto.

\begin{quadro}[htb]
\centering
\caption{Roteiro de tarefas do estudo empírico}
\label{quad:tarefas}
\footnotesize
\begin{tabular}{c p{4.4cm} p{3.4cm} p{3.2cm}}
\toprule
\textbf{\#} & \textbf{Cenário apresentado} & \textbf{Elementos verificados} &
\textbf{Critério de conclusão} \\
\midrule
T1 & Descobrir o que o aplicativo faz e escolher uma trilha & Tela inicial e
listagem; compreensão da proposta & Alcança a tela de uma trilha e verbaliza a
proposta \\
\addlinespace
T2 & Relatar o que a trilha escolhida exigirá & Tela de preparação; lista de
verificação & Verbaliza distância, número de \textit{checkpoints} e exigências \\
\addlinespace
T3 & Chegar ao primeiro \textit{checkpoint} & Mapa, marcadores, bússola,
permissão de localização & Chega ao ponto sem auxílio \\
\addlinespace
T4 & Registrar a chegada ao \textit{checkpoint} & Leitura de QR Code; permissão de
câmera; estados de retorno & Realiza a leitura e reconhece a confirmação \\
\addlinespace
T5 & Relatar o estado corrente da atividade sem tocar na tela & Progresso,
cronômetro, distância & Interpreta corretamente os três indicadores \\
\addlinespace
T6 & Retomar a atividade após interrupção externa & Suspensão e retomada & Retoma
no ponto correto sem reiniciar \\
\addlinespace
T7 & Concluir a trilha e relatar o desempenho & Tela de conclusão & Interpreta
tempo e distância finais \\
\bottomrule
\end{tabular}
\\[4pt]
\begin{flushleft}
Fonte: Elaborado pelo autor.
\end{flushleft}
\end{quadro}

Registra-se ainda uma observação não induzida: caso o participante alcance
espontaneamente a tela de perfil, cuja limitação está registrada no
Quadro~\ref{quad:mapeamento}, sua reação é integralmente documentada. A ausência de
visita espontânea constitui, igualmente, dado.

\section{Procedimento de Coleta}
\label{sec:procedimento}

Cada sessão segue roteiro padronizado, idêntico para todos os participantes, com
duração aproximada de setenta e cinco minutos, distribuída conforme o
Quadro~\ref{quad:sessao}.

\begin{quadro}[htb]
\centering
\caption{Estrutura da sessão de avaliação}
\label{quad:sessao}
\begin{tabular}{p{6.4cm} p{2.0cm} p{4.4cm}}
\toprule
\textbf{Etapa} & \textbf{Duração} & \textbf{Registro} \\
\midrule
Acolhimento e enquadramento & 5 min & Áudio \\
Leitura e assinatura do TCLE & 5 min & Documento \\
Apresentação mínima do aplicativo & 3 min & Áudio \\
Explicação do protocolo & 4 min & Áudio \\
Treino de verbalização & 5 min & Áudio \\
Coleta de dados demográficos & 3 min & Ficha \\
Execução das tarefas & 30--45 min & Áudio, tela, ficha \\
Recuperação retrospectiva assistida & 5 min & Áudio \\
Item de esforço percebido & 1 min & Ficha \\
Entrevista semiestruturada & 10--15 min & Áudio \\
Instrumento TAM & 7 min & Formulário \\
\bottomrule
\end{tabular}
\\[4pt]
\begin{flushleft}
Fonte: Elaborado pelo autor.
\end{flushleft}
\end{quadro}

A sessão é conduzida por dois pesquisadores, com papéis separados: um modera e
cronometra, outro preenche a ficha de observação e registra os instantes dos
episódios de atrito, posteriormente utilizados para selecionar os trechos da
recuperação retrospectiva assistida. O roteiro integral de moderação, incluindo as
falas de enquadramento, as tarefas e as regras de intervenção, consta do
Apêndice~\ref{ape:roteiro}.

\subsection{Ficha de observação}

Cada episódio de atrito é registrado em uma linha da ficha estruturada, com os
campos: tarefa; instante de ocorrência; conclusão da tarefa, classificada em
autônoma, com auxílio --- indicando o grau da ajuda escalonada --- ou falha;
verbalização literal associada; heurística tocada, conforme a numeração do
Quadro~\ref{quad:heuristicas}; origem do atrito, conforme a classificação da
Seção~\ref{sec:gamificacao}; e severidade, conforme o
Quadro~\ref{quad:severidade}.

\section{Instrumento de Aceitação Tecnológica}
\label{sec:instrumento_tam}

\subsection{Composição do instrumento}

O instrumento foi elaborado a partir dos construtos originais de
\citeonline{davis1989}, com adaptação da redação dos itens ao contexto de uso do
aplicativo. É composto por dez itens, distribuídos em três construtos, e aplicado
em escala Likert de sete pontos ancorada nos extremos por ``Discordo Totalmente''
e ``Concordo Totalmente'', com ponto neutro central. O instrumento é
autoaplicado, em formulário eletrônico, imediatamente após a entrevista, ainda em
campo. O Quadro~\ref{quad:tam} apresenta os itens conforme aplicados.

\begin{quadro}[htb]
\centering
\caption{Itens do instrumento baseado no Technology Acceptance Model}
\label{quad:tam}
\footnotesize
\begin{tabular}{p{1.4cm} p{11.4cm}}
\toprule
\textbf{Código} & \textbf{Item} \\
\midrule
\multicolumn{2}{l}{\textbf{Facilidade de Uso Percebida (PEOU)}} \\
\midrule
PEOU1 & Minha interação com o aplicativo \textit{Trilha das Árvores} é clara e
compreensível \\
PEOU2 & Interagir com o aplicativo \textit{Trilha das Árvores} não exige muito do
meu esforço mental \\
PEOU3 & Considero o aplicativo \textit{Trilha das Árvores} fácil de usar \\
PEOU4 & Eu acho fácil conseguir que o aplicativo \textit{Trilha das Árvores} faça
o que eu quero que ele faça \\
\midrule
\multicolumn{2}{l}{\textbf{Utilidade Percebida (PU)}} \\
\midrule
PU1 & Usar o aplicativo \textit{Trilha das Árvores} melhorou o meu desempenho na
realização da trilha\footnotemark \\
PU2 & Usar o aplicativo \textit{Trilha das Árvores} na minha visita melhorou a
minha produtividade na exploração das trilhas \\
PU3 & Usar o aplicativo \textit{Trilha das Árvores} aumenta a minha eficácia na
identificação das árvores \\
PU4 & Eu considero o aplicativo \textit{Trilha das Árvores} útil para explorar o
parque \\
\midrule
\multicolumn{2}{l}{\textbf{Intenção de Uso (ITU)}} \\
\midrule
ITU1 & Supondo que eu tenha acesso ao aplicativo \textit{Trilha das Árvores}, eu
pretendo usá-lo \\
ITU2 & Levando em conta que eu tenho acesso ao aplicativo \textit{Trilha das
Árvores}, eu prevejo que eu vou usá-lo \\
\bottomrule
\end{tabular}
\\[4pt]
\begin{flushleft}
Fonte: Elaborado pelo autor, com base em \citeonline{davis1989}.
\end{flushleft}
\end{quadro}

\footnotetext{O item PU1 é acompanhado, no formulário, de uma âncora de comparação
apresentada ao respondente, que o instrui a considerar como referência a
realização do mesmo percurso com um mapa físico, sem os recursos do aplicativo.}

\subsection{Justificativa da âncora de comparação}

O construto de Utilidade Percebida, tal como formulado por
\citeonline{davis1989}, pressupõe uma referência implícita: o desempenho
\textit{sem} o sistema. Em contextos organizacionais, essa referência é evidente
para o respondente, que conhece o processo anterior à informatização. No contexto
deste estudo, o participante realiza a trilha \textbf{apenas} com o aplicativo e
não dispõe de experiência prévia do percurso sem mediação, o que tornaria a
comparação indeterminada. A explicitação da âncora --- o percurso com mapa físico
--- restitui ao item a referência que o construto pressupõe. Por essa razão, a
comparação com o mapa físico é deliberadamente \textbf{excluída} do roteiro da
entrevista que precede o instrumento, a fim de evitar efeito de priorização sobre
a resposta.

\subsection{Função do instrumento na primeira rodada}

Com três participantes, os escores obtidos na Fase 1 não admitem tratamento
inferencial nem estimativa confiável de consistência interna. O instrumento cumpre,
nessa fase, função de \textbf{validação semântica}: verifica-se se cada item é
compreendido de forma unívoca, solicitando ao participante que leia o enunciado em
voz alta e relate o que dele entendeu antes de assinalar a resposta. Registram-se
as reformulações espontâneas e os termos que suscitem dúvida. Os escores são
reportados de forma descritiva e individual, sem agregação.

% ============================================================================
% >>> [PREENCHER --- AJUSTES DE INSTRUMENTAÇÃO APLICADOS] <<<
% Registrar aqui, para efeito de rastreabilidade, os ajustes realizados no
% formulário antes da aplicação. Pontos já identificados:
%
%   (a) Anonimato: o formulário coletava o endereço de correio eletrônico do
%       respondente, o que contradiz a declaração de anonimato do próprio
%       cabeçalho e do TCLE. Desativar a coleta e substituir por código
%       sequencial de participante (P1, P2, P3).
%   (b) Redação da âncora do item PU1: revisar a frase de apoio, cuja
%       formulação atual apresenta erro de concordância e ambiguidade.
%   (c) Descrição do bloco de Intenção de Uso: o texto de apoio menciona
%       recomendação a terceiros, dimensão não contemplada por nenhum dos dois
%       itens. Ajustar a descrição ou acrescentar item correspondente.
%   (d) Escala: confirmada em sete pontos nos dez itens.
% ============================================================================

\section{Plano de Análise}
\label{sec:analise}

\subsection{Dados qualitativos}

As gravações de áudio são integralmente transcritas. As transcrições, articuladas
aos registros de tela e às fichas de observação, são submetidas a análise temática
\cite{brauneclarke2006}, com codificação de cada episódio de atrito segundo dois
eixos independentes: a heurística violada e a origem do atrito. A codificação
inicial é revista em segunda passagem, e os casos de classificação ambígua são
registrados como tal, com a justificativa da decisão adotada.

\subsection{Dados quantitativos}

Os dados do instrumento TAM são reportados de forma descritiva, por construto e
por participante. Os tempos de tarefa, as taxas de conclusão e os graus de auxílio
são igualmente reportados de forma descritiva. Em rodadas futuras, com amostra
ampliada, prevê-se a verificação da consistência interna dos construtos por
coeficiente alfa de \citeonline{cronbach1951}.

\subsection{Triangulação e matriz de rastreabilidade}

A etapa final consiste na construção da matriz de rastreabilidade, instrumento que
materializa a contribuição enunciada na Seção~\ref{sec:problema}. Cada linha da
matriz associa um problema observado às seguintes dimensões: a heurística violada;
a decisão de engenharia que o originou, conforme o Quadro~\ref{quad:mapeamento}; a
origem do atrito; a severidade; o construto do TAM potencialmente impactado; e a
correção proposta. O Quadro~\ref{quad:matriz_modelo} apresenta a estrutura da
matriz, cujo preenchimento consta do Capítulo~\ref{cap:resultados}.

\begin{quadro}[htb]
\centering
\caption{Estrutura da matriz de rastreabilidade}
\label{quad:matriz_modelo}
\footnotesize
\begin{tabular}{p{1.0cm} p{2.6cm} p{2.6cm} p{1.6cm} p{1.4cm} p{2.4cm}}
\toprule
\textbf{ID} & \textbf{Problema observado} & \textbf{Decisão de origem} &
\textbf{Origem} & \textbf{Sev.} & \textbf{Correção proposta} \\
\midrule
P01 & --- & --- & --- & --- & --- \\
P02 & --- & --- & --- & --- & --- \\
\bottomrule
\end{tabular}
\\[4pt]
\begin{flushleft}
Fonte: Elaborado pelo autor.
\end{flushleft}
\end{quadro}

\section{Ameaças à Validade}
\label{sec:ameacas}

Registram-se as principais ameaças à validade do estudo e as medidas de mitigação
adotadas.

\begin{alineas}
  \item \textbf{Reatividade da verbalização.} A exigência de falar durante a
  tarefa pode alterar o desempenho. Mitigação: protocolo híbrido segmentado, com
  verbalização concomitante restrita aos segmentos estacionários, e treino prévio.

  \item \textbf{Viés do moderador.} O moderador é também o desenvolvedor do
  sistema, o que introduz risco de indução involuntária. Mitigação: roteiro
  fechado de sondas, lista explícita de formulações proibidas, separação de papéis
  entre moderador e anotador, e enquadramento inicial que autoriza a crítica.

  \item \textbf{Desejabilidade social.} O participante pode moderar críticas por
  cortesia. Mitigação: enquadramento explícito de que o objeto avaliado é o
  sistema; anonimato assegurado; triangulação entre o observado e o declarado.

  \item \textbf{Dimensão da amostra.} Três participantes por rodada não sustentam
  generalização. Mitigação: caráter declaradamente formativo do estudo e ausência
  de pretensão inferencial, conforme a Seção~\ref{sec:escopo}.

  \item \textbf{Variação ambiental entre rodadas.} Diferenças de clima ou de
  horário poderiam ser confundidas com efeito do redesenho. Mitigação:
  padronização e registro das condições, com replicação na Fase 3.

  \item \textbf{Efeito de aprendizagem.} Mitigação: participantes distintos entre
  rodadas, conforme a Seção~\ref{sec:participantes}.
\end{alineas}

\section{Considerações Finais}
\label{sec:met_finais}

Este capítulo detalhou o desenho da avaliação em suas quatro fases, a
caracterização do público por proto-personas, os critérios de recrutamento e os
aspectos éticos, o aparato técnico e a estratégia de padronização das condições, o
roteiro de tarefas, o procedimento de coleta, o instrumento de aceitação
tecnológica e o plano de análise, encerrando com o registro das ameaças à validade
e das medidas de mitigação. O capítulo seguinte apresenta os resultados obtidos.
